\section{Introduction}
\label{sec:introduction}
Social media is full of scientific claims, but they are usually oversimplified, informal, and stripped of context. For fact-checkers, tracking down the original study behind a tweet or post is a painfully slow task—especially when the post paraphrases the study or uses a different language. Task 1 of the CLEF 2026 CheckThat! Lab~\cite{clef-checkthat:2026:task1} tackles this issue. The goal is to automatically link social media claims to their scientific sources. This comes with two major hurdles: a massive style gap between casual internet slang and formal abstracts, and a cross-lingual challenge (matching English, French, and German posts against a mostly English corpus).We address this with a highly configurable IR system. First, we translate all non-English text into English to keep our pipeline consistent. The retrieval strategy pairs traditional BM25 search (utilizing structured fields like title, abstract, and authors) with dense multilingual embeddings to capture semantic similarity. Our system supports a wide range of preprocessing and indexing configurations, allowing us to experiment with different setups to bridge the vocabulary gap. Finally, we evaluate the impact of query expansion and cross-encoder reranking on overall retrieval performance.

This paper is organized into five sections. Section~\ref{sec:related-work} reviews some related work; Section~\ref{sec:methodology} presents our approach in detail; Section~\ref{sec:setup} describes the experimental setup; Section~\ref{sec:results} discusses our main findings and analyzes the obtained results; and Section~\ref{sec:conclusion} concludes the paper and describes possible future work.